\chapter{Prompts e Esquemas do Oráculo LLM}
\label{ap:prompts}

\section{Prompt v3 (produção)}

Contexto de tarefa comum a todos os modos (o texto integral e versionado
acompanha o repositório de código da tese):

\begin{quote}\small
``Você é um especialista em catalogação de produtos do varejo brasileiro. As
descrições vêm de cupons fiscais e cadastros de e-commerce: são curtas,
geralmente em caixa alta e cheias de abreviações (ex.: CERV=cerveja,
REFR=refrigerante, BISC=biscoito, CHOC=chocolate, RECH=recheado, PCT=pacote,
LT=lata, CX=caixa, UN=unidade). Expanda mentalmente as abreviações e
classifique o TIPO do produto, ignorando marca, tamanho e embalagem. Escolha
sempre a categoria mais específica aplicável; na dúvida entre uma categoria
plausível e `\_rare\_', prefira a categoria plausível — use `\_rare\_' SOMENTE
se nenhuma categoria da lista descrever o tipo do produto. Exemplos:
`CERV BRAHMA LT 350ML' $\to$ cerveja; `BISC RECH CHOC 130G' $\to$ biscoito;
`REQUEIJAO LIGHT CP 200G' $\to$ requeijao.''
\end{quote}

O prefixo estático (contexto + esquema) precede o item variável, habilitando o
\textit{cache} de prefixo dos provedores (88--95\% de acerto medido no E0).

\section{Modos de instrumentação}

\textbf{enum}: saída estruturada com \texttt{strict:true}; o campo
\texttt{predicted\_category} é restrito por \texttt{enum} às 621 categorias do
\textit{CategorySchema}; campos \texttt{expanded\_description} e
\texttt{rationale} completam o objeto (todos obrigatórios, exigência do modo
estrito). \textbf{json-prompt} (provedores sem saída estruturada): a lista de
categorias e o formato JSON são instruídos no texto; a resposta passa por
extração tolerante e validação exata contra o esquema, com rótulos fora do
espaço contabilizados como inválidos. \textbf{free} (somente RQ4): saída
estruturada sem \texttt{enum}: réplica controlada do instrumento defeituoso
do trabalho anterior. Na rotulagem em lote, os itens são numerados e a
resposta é um vetor indexado; o esquema do lote não fixa
\texttt{minItems}/\texttt{maxItems} para preservar a cacheabilidade do
prefixo; o uso (tokens/custo) é rateado igualmente entre os itens da chamada.

\section{Variantes do E0-P}

\textbf{v4a} acrescenta dez regras de fronteira do catálogo derivadas da
anatomia de erros (ex.: medicamentos $\to$ \textit{outro farma}; produto
bucal líquido $\to$ \textit{antisseptico bucal}; pó para suco $\to$
\textit{preparo para suco}; fones com fio $\to$ \textit{acessorio de audio});
\textbf{v4b} soma dez exemplos \emph{inventados} de fronteiras difíceis
(nunca instâncias das amostras de avaliação, decisão D-004). Texto integral
das variantes no mesmo módulo versionado.
